\section{Alertas personales persistentes de actividad}\label{sec:implementacion-alertas}

Las alertas de audiencias y plazos constituyen un servicio distinto de la consulta de agenda. Conservan preferencias personales, planificación, avisos internos, lectura e intentos de correo. Sus cuatro motivos son proximidad, vencimiento sin atención declarada, revisión humana requerida y cambio de fecha próximo. Las anticipaciones iniciales son 48 y 24 horas; cada familia admite hasta ocho valores distintos entre 1 y 720 horas, con canales interno y correo independientes. Estos intervalos expresan tiempo transcurrido respecto de un instante ya determinado, sin introducir otro cómputo jurídico. La ventana de proximidad incluye su umbral y excluye el vencimiento, con precisión de nanosegundos. Al recuperar actividad después de una interrupción, se selecciona el umbral cruzado más próximo que continúa siendo elegible.

La cuenta autenticada administra únicamente sus preferencias y su bandeja. Las audiencias se dirigen a miembros activos del expediente con un rol interno autorizado; Owner sólo es destinatario cuando pertenece al expediente. Un plazo se dirige a su responsable vigente si conserva autorización. Las fichas de participantes procesales no se convierten en cuentas receptoras y Client permanece denegado. La consulta vuelve a comprobar destinatario y acceso actual antes de aplicar el límite. Un cierre administrativo no suspende términos ni elimina por sí solo los avisos. Antes de activar un aviso o autorizar un intento de correo se verifica nuevamente el recurso: una fecha histórica no sustituye el vencimiento operativo ni la comprobación actual de dependencias descrita en la \cref{sec:implementacion-plazos-persistentes}.

\subsection{Planes, episodios y confirmación auditada}

Las migraciones del grupo \texttt{0021} incorporan ocho tablas PostgreSQL para preferencias, estado del recurso, cursor de recorrido, planes, notificaciones, recibos de lectura, entregas e intentos. Las preferencias conservan revisión y operación idempotente; una repetición exacta recupera su resultado original y una operación reutilizada con otros valores produce conflicto. La revisión cero representa los valores iniciales sin inventar recibo ni fecha. El arranque comprueba catálogo, funciones, restricciones, disparadores, propietarios, permisos e inventario. El rol operativo recibe lectura, inserción por columnas y actualizaciones acotadas; carece de facultades para borrar evidencia o reinsertar el cursor. Las preferencias, recibos de lectura e intentos conservan filas inmutables.

El escáner mantiene progreso durable y procesa cantidades acotadas de recursos y destinatarios. Los escritores de audiencias y plazos invalidan su estado de alertas dentro de la misma transacción auditada. La reconciliación verifica el origen y las preferencias; la activación confirma conjuntamente el aviso, la entrega pendiente cuando corresponde y su auditoría. Los errores de escritura revierten esas modificaciones. Una ocurrencia de proximidad se vincula al recurso, destinatario, instante exacto y anticipación; una nueva revisión sin cambio temporal no obliga a repetirla. Un plan sustituido que nunca produjo una notificación puede recuperar elegibilidad, mientras una ocurrencia ya activada conserva su identidad.

La falta de atención y la revisión requerida se organizan en episodios. Registrar atención o aceptar una corrección humana cierra el episodio correspondiente dentro de la escritura del plazo. Esta transición se conserva aunque se reabra la atención o cambie otra fuente antes del siguiente escaneo. Al consultar el aviso anterior se verifica la revisión histórica que terminó su episodio; el nuevo estado puede producir otra ocurrencia sin reactivar la anterior. Leer un aviso no registra atención, no acepta la aplicabilidad del perfil y no modifica la historia del recurso.

\subsection{Correo, composición y presentación}

La entrega externa se separa de la transacción PostgreSQL. Un arrendamiento identifica el intento vigente y permite rechazar confirmaciones atrasadas; los intentos y su proyección se conservan con auditoría. El primer intento fija remitente, destinatario, plantilla, URL y clave idempotente. El adaptador seleccionado es Resend y los reintentos conservan exactamente ese mensaje. La implementación limita los intentos a ocho y la ventana de reintento a 23 horas; una respuesta incierta fuera de esos límites requiere conciliación y no recibe automáticamente otra clave. El estado aceptado significa aceptación por el proveedor, sin acreditar entrega al buzón ni lectura. El correo contiene sólo un aviso genérico y un enlace fijo de acceso, sin datos del expediente. Una revocación impide autorizar un intento posterior, pero no puede retirar una petición que ya salió hacia el proveedor.

El ejecutable \texttt{serve} comparte el almacén PostgreSQL entre el servicio HTTP, la generación periódica y el consumidor de correo. El transporte externo exige configuración completa de credencial, remitente y URL; su ausencia conserva el canal deshabilitado y permite la bandeja interna. Los ciclos tienen presupuesto y pausa configurables, con 20 acciones y 1000 milisegundos como valores iniciales. La supervisión coordina HTTP, alertas y reevaluación de plazos, espera las operaciones bloqueantes en curso y libera los adaptadores al terminar. La composición está implementada y tiene aceptación con HTTP, restauración y navegador reales, con el alcance de la \cref{sec:pruebas-alertas}.

Qadra presenta preferencias y bandeja con filtros de lectura y resolución, continuación y estados de correo separados. Conserva el comando ante una respuesta incierta y permite contrastar su recibo antes de repetirlo. Abrir un aviso revalida el expediente y recupera exactamente la revisión de origen; no la sustituye por la cabeza actual. El título y la referencia proceden del contexto histórico capturado. La alerta no acredita vigencia futura, actuación oportuna ni calificación jurídica. Tampoco modifica los límites de la CA interna y de la TSA local: sus evidencias no constituyen FIREL ni una constancia de un PSC autorizado. El contrato y la decisión están en \rutarepo{docs/alerts-api.md} y \rutarepo{docs/adr/0039-durable-activity-alerts.md}.
